\subsection{性、爱与承诺的关系}

性、爱与承诺是亲密关系的三个组成部分，相互影响、互为条件。三者的组合方式决定关系的形态：有性无爱可以是双方同意的明确约定；有爱无性可能是阶段性的（疾病、产后、压力期），也可能是长期的生活方式；有承诺而另两者稀薄的关系，则容易走向"室友化"。真正需要关注的，不是三者是否存在，而是它们之间的落差如何被处理——双方能否就"我们现在的样子"达成诚实共识，而不是各自按期待默默扣分。

\begin{itemize}
  \item \keyword{性与爱}：性可以促进爱，爱可以提升性；性爱分离（如一夜情、性交易）虽存在，但通常难以带来持久的满足。
  \item \keyword{性与承诺}：承诺为性提供安全感与约束力，性也可以成为表达承诺、巩固信任的方式。
  \item \keyword{三者的平衡}：培养基于爱的性、以承诺为性提供稳定的环境，并保持三者协调发展，避免任何一方的长期缺失。
\end{itemize}

一个可操作的练习：双方各自为"性、爱、承诺"三项打 0--10 分，写下希望一年后变化的方向，以及自己愿意为此做的一件具体的事；随后交换、不打断地听完，并复述对方的答案以确认理解。这一练习通常不产出"解决方案"，而是把模糊的委屈转化为可以讨论的具体议题。
